\sgaStableTarget{sga6.e06.cor.6-3}\textbf{Corollary 6.3.} \uline{With the hypotheses of} 6.1, \uline{let} @@M0@@ \uline{be a commutative ring}. \uline{For an element} @@M1@@ \uline{of} @@M2@@ \uline{to be invertible}, \uline{it is necessary and sufficient that} @@M3@@ \uline{be invertible}. \uline{In particular}, \uline{if} @@M4@@ \uline{is a locally free sheaf of finite type on} @@M5@@, \uline{everywhere} @@M6@@, \uline{then its class in} @@M7@@ \uline{is invertible}.

The augmentation ideal of @@M8@@ is @@M9@@, and is therefore a nil ideal by 6.1. Thus, for an element of @@M10@@ to be invertible, it is necessary and sufficient that its augmentation be invertible.

\sgaStableTarget{sga6.e06.cor.6-4}\textbf{Corollary 6.4.} \uline{Let} @@M11@@ \uline{be a quasi}-\uline{compact scheme}, \uline{having an ample sheaf}, \uline{and} @@M12@@ \uline{a proper and perfect morphism}. \uline{Suppose in addition that one of the following two conditions is satisfied}:
@@M13@@
\uline{so that the} ``\uline{finiteness theorem}'' \uline{holds} (III 2.2). \uline{We are given a ring} @@M14@@, \uline{and} \uline{consider the} ``\uline{direct image}'' \uline{homomorphism} (IV 2.11):
@@M15@@
\uline{Finally suppose that there is an element} @@M16@@ \uline{such that} @@M17@@ \uline{is invertible}. \uline{Then the homomorphism}:
@@M18@@
\uline{is injective}.

We write the projection formula:
@@M19@@
Now @@M20@@ has invertible augmentation, so it is invertible (6.3). Consequently, the composite homomorphism:
